\documentclass[a4paper,10pt,fleqn]{article}
\usepackage[margin=1in]{geometry}
\usepackage{amssymb}
\usepackage{adjustbox}
\usepackage{natbib}
\usepackage[colorlinks=true,linkcolor=blue,citecolor=blue,urlcolor=blue]{hyperref}
\usepackage{fuzz}
\usepackage{zed-maths}
\usepackage{zed-proof}
\newdimen\savedleftskip
\title{Primary Keys in Relational Schemas}
\author{txt2tex example}
\hypersetup{pdftitle={Primary Keys in Relational Schemas},pdfauthor={txt2tex example},pdfcreator={txt2tex v1.2.0}}
\begin{document}

\maketitle

\begin{zed}[ClassName, ShipName, BattleName, CountryName, PersonID, CentreID]\end{zed}

\section*{Single Primary Key}

\begin{schema}{Class}
class : ClassName \\
country : CountryName \\
numGuns : \nat \\
bore : \nat \\
displacement : \nat
\end{schema}
\noindent$\mathrm{PK}(\mathrm{Class}) = \{class\}$

\begin{schema}{Ship}
name : ShipName \\
class : ClassName \\
launched : \nat
\end{schema}
\noindent$\mathrm{PK}(\mathrm{Ship}) = \{name\}$

\begin{schema}{Battle}
name : BattleName \\
date : \nat
\end{schema}
\noindent$\mathrm{PK}(\mathrm{Battle}) = \{name\}$

\section*{Composite Primary Key}

\begin{schema}{CentreStaff}
centreId : CentreID \\
staffId : PersonID
\end{schema}
\noindent$\mathrm{PK}(\mathrm{CentreStaff}) = \{centreId, staffId\}$

\end{document}